\uuid{7Mva}
\exo7id{1216}
\titre{exo7 1216}
\auteur{ridde}
\organisation{exo7}
\datecreate{1999-11-01}
\isIndication{false}
\isCorrection{true}
\chapitre{Continuité, limite et étude de fonctions réelles}
\sousChapitre{Fonctions équivalentes, fonctions négligeables}
\cpge{mpsi-mp2i.limites-continuite.a, pcsi.limites-continuite.a, ptsi.limites-continuite.a, mpsi-mp2i.limites-continuite.c, pcsi.limites-continuite.c, ptsi.limites-continuite.c, bcpst1.limites-continuite.a, bcpst1.limites-continuite.b, bcpst1.developpements-limites.b}
\module{Analyse}
\niveau{L1}
\difficulte{}

\contenu{
\question{
\'Etudier en $ + \infty$ et $-\infty$ la fonction
$f (x) = \sqrt[3]{x^3 + 1}
 + \sqrt{x^2 + x + 1}$.
}
\reponse{
$$\lim _{x\rightarrow \infty }\sqrt [3]{{x}^{3}+1}-\sqrt {{x}^{2}+x+1}=-
1/2.$$
}
}
